\chapter{Optimal Stopping and Impulse Control}\label{ch:qm:optimal-stopping-and-impulse-control}

A bank's treasury desk carries a euro exposure that drifts with its clients' flows: every payment,
every trade finance deal, moves it by a few million. Each hedge ticket costs a fixed amount; carrying the
exposure costs a risk charge that grows with its square. Hedging after every client flow wastes fees;
never hedging lets the risk grow without bound. The answer is a band: do nothing while the exposure is
inside $\pm b$, trade it back to zero when it touches the edge. With the desk's numbers the best band is
$\pm$USD~34.6 million, a hedge every three days on average, and the policy costs half as much as hedging
once a day at the close; if the ticket fee halves, the band narrows by only 16\%, because the optimal band
grows with the fourth root of the fee. Deciding when to act is a stopping problem; deciding when and how
much to act, repeatedly, is \omterm{def:qm:optimal-stopping-and-impulse-control:impulse}{impulse control}. This chapter solves both: the \omterm{def:qm:optimal-stopping-and-impulse-control:snell}{Snell envelope} in discrete time,
free boundaries and \omterm{def:qm:optimal-stopping-and-impulse-control:fbp}{smooth pasting} in continuous time, \omterm{def:qm:optimal-stopping-and-impulse-control:impulse}{impulse control} and its \omterm{def:qm:optimal-stopping-and-impulse-control:notrade}{no-trade region}, and
\omterm{def:qm:optimal-stopping-and-impulse-control:singular}{singular control}, its limit when costs are proportional.

\section{Optimal stopping in discrete time: the Snell envelope}

\begin{definition}[Optimal stopping problem, Snell envelope]\label{def:qm:optimal-stopping-and-impulse-control:snell}
Given an adapted, integrable reward process $(G_n)_{n\le N}$, the \emph{optimal stopping problem}\index{optimal
stopping problem} is to find $\sup_\tau\E[G_\tau]$ over \omterm{def:qm:probability-at-speed:stopping}{stopping times} $\tau \le N$. Its \emph{Snell envelope}\index{Snell
envelope} is defined backward by $U_N = G_N$ and $U_n = \max(G_n, \E_n[U_{n+1}])$.
\end{definition}

\begin{theorem}[The Snell envelope solves the stopping problem]\label{thm:qm:optimal-stopping-and-impulse-control:snell}
$U$ is the smallest \omterm{def:qm:probability-at-speed:martingale}{supermartingale} dominating $G$; $U_n = \sup_{\tau \ge n}\E_n[G_\tau]$; and $\tau^* = \min\{n : U_n = G_n\}$ is
optimal.
\end{theorem}

\begin{proof}
$U_n \ge \E_n[U_{n+1}]$ and $U_n \ge G_n$ by construction. If $Y$ is a \omterm{def:qm:probability-at-speed:martingale}{supermartingale} dominating $G$, then $Y_N \ge
U_N$ and, backward, $Y_n \ge \max(G_n, \E_n[Y_{n+1}]) \ge U_n$. Up to $\tau^*$ the envelope equals its \omterm{def:qm:probability-at-speed:condexp}{conditional
expectation}, so $U_{n\wedge\tau^*}$ is a \omterm{def:qm:probability-at-speed:martingale}{martingale} and $\E[G_{\tau^*}] = \E[U_{\tau^*}] = U_0$; any other $\tau$ gives $\E[G_\tau] \le
\E[U_\tau] \le U_0$ by optional stopping for the \omterm{def:qm:probability-at-speed:martingale}{supermartingale}.
\end{proof}

\begin{definition}[Continuation region, stopping region]\label{def:qm:optimal-stopping-and-impulse-control:regions}
For a Markov problem with $G_n = g(n, X_n)$ and $U_n = u(n, X_n)$, the \emph{stopping region}\index{stopping region} is
$\{(n, x) : u = g\}$ and the \emph{continuation region}\index{continuation region} its complement, where waiting is
strictly better.
\end{definition}

The running project's solver of chapter~9 computes the \omterm{def:qm:optimal-stopping-and-impulse-control:snell}{Snell envelope} by adding the stopping reward to the
backward induction: $V_t = \max(\text{stop}, \text{continue})$. An American option is the best-known stopping
problem; One Quant Book~5, chapter~6, prices it and names its boundary. Here the same mathematics decides
when to take profit on a trade.

\section{Free boundaries and variational inequalities}

In continuous time, with $X$ a diffusion with generator $\mathcal L$ and a discount rate $r$, the value $V(x) =
\sup_\tau\E_x[e^{-r\tau}g(X_\tau)]$ of a perpetual stopping problem satisfies, where smooth,
\[
\max\bigl(\mathcal LV - rV,\ g - V\bigr) = 0:
\]
in the \omterm{def:qm:optimal-stopping-and-impulse-control:regions}{continuation region} $V > g$ and $V$ solves the equation $\mathcal LV = rV$; in the \omterm{def:qm:optimal-stopping-and-impulse-control:regions}{stopping region} $V = g$ and
waiting would not pay, $\mathcal Lg - rg \le 0$.

\begin{definition}[Variational inequality, free-boundary problem, smooth pasting]\label{def:qm:optimal-stopping-and-impulse-control:fbp}
The display above is a \emph{variational inequality}\index{variational inequality}. Written as ``solve $\mathcal LV = rV$ on
an unknown region $D$ with $V = g$ on its boundary'', it is a \emph{free-boundary problem}\index{free-boundary problem}:
the boundary $\partial D$ is part of the solution. The extra condition that fixes it, $V' = g'$ on $\partial D$ where $g$ is
smooth, is \emph{smooth pasting}\index{smooth pasting} (or smooth fit).
\end{definition}

\begin{proposition}[Taking profit on a drifting position]\label{prop:qm:optimal-stopping-and-impulse-control:drift}
Let $X_t = x + \mu t + \sigma W_t$ be the running P\&L of a position, $c$ the cost of closing it, and $r > 0$ a discount rate.
The optimal rule for $\sup_\tau\E[e^{-r\tau}(X_\tau - c)]$ closes the first time $X$ reaches $b^* = c + 1/\theta$, where $\theta > 0$ solves
$\tfrac12\sigma^2\theta^2 + \mu\theta - r = 0$, and $V(x) = (b^* - c)e^{\theta(x - b^*)}$ for $x < b^*$.
\end{proposition}

\begin{proof}
For a threshold $b$, $\E_x[e^{-r\tau_b}] = e^{\theta(x - b)}$ for $x < b$, since $e^{-rt + \theta X_t}$ is a \omterm{def:qm:probability-at-speed:martingale}{martingale} exactly when
$\theta$ solves the quadratic; so the value of the $b$-rule is $(b - c)e^{\theta(x - b)}$, maximised over $b$ at $b = c + 1/\theta$.
There $V'(b^*) = \theta(b^* - c) = 1 = g'(b^*)$: \omterm{def:qm:optimal-stopping-and-impulse-control:fbp}{smooth pasting}. On $x \ge b^*$, $V = x - c$ and $\mathcal Lg - rg = \mu - r(x - c) \le 0$
because $x - c \ge 1/\theta \ge \mu/r$, so the \omterm{def:qm:optimal-stopping-and-impulse-control:fbp}{variational inequality} holds and a verification argument as in chapter~9
applies.
\end{proof}

With $\mu = 0.5$ and $\sigma = 2$ basis points a day, a closing cost of 5 bp and $r = 2\%$ a day, $\theta = 0.0351$ and $b^* = 33.5$ bp.
Backward induction on a grid, observing the position once a day, stops at 32.2 bp and values the position at 8.79
bp at zero against 8.80 (\cref{fig:qm:optimal-stopping-and-impulse-control:stopping}): the daily grid stops a
little earlier, by about the overshoot $0.5826\,\sigma\sqrt{\Delta t} = 1.2$ bp of chapter~2.

\begin{omfigure}
\begin{tikzpicture}
\begin{axis}[width=11cm,height=4.8cm,
  xlabel={running P\&L of the position $x$ (bp)},ylabel={value (bp)},
  tick label style={font=\small},label style={font=\small},
  xmin=-10,xmax=60,ymin=0,ymax=56,grid=major,grid style={gray!25},
  legend columns=2,legend style={font=\footnotesize,at={(0.5,-0.3)},anchor=north,draw=none,fill=none,/tikz/every even column/.append style={column sep=8pt}},
  table/col sep=comma]
\addplot[omDef,very thick] table[x=x,y=value]{figdata/methods/10-optimal-stopping-and-impulse-control/stopping.csv};
\addplot[omThm,thick,dashed] table[x=x,y=payoff]{figdata/methods/10-optimal-stopping-and-impulse-control/stopping.csv};
\addplot[gray,thin,forget plot] coordinates {(33.51,0) (33.51,56)};
\node[font=\footnotesize,anchor=north west,fill=white,inner sep=1pt] at (axis cs:34.5,52) {$b^* = 33.5$};
\legend{value $V(x)$,{close now: $x - c$}}
\end{axis}
\end{tikzpicture}

\omcaption{Taking profit on a drifting position ($\mu = 0.5$, $\sigma = 2$ bp a day, closing cost 5 bp, discount 2\% a day). The
value of waiting (solid) lies above the value of closing now (dashed) until $b^* = c + 1/\theta$, where it meets it with the same
slope: \omterm{def:qm:optimal-stopping-and-impulse-control:fbp}{smooth pasting}. Data: closed form, checked by the chapter's grid
solver.}\label{fig:qm:optimal-stopping-and-impulse-control:stopping}
\end{omfigure}

\section{Impulse control and no-trade regions}

\begin{definition}[Impulse control, quasi-variational inequality]\label{def:qm:optimal-stopping-and-impulse-control:impulse}
\emph{Impulse control}\index{impulse control} acts on a state by discrete interventions: at \omterm{def:qm:probability-at-speed:stopping}{stopping times} $\tau_1 < \tau_2 <
\dots$ it moves $X$ by jumps $\xi_i$, each paying a cost with a fixed part $K > 0$. Its \omterm{def:qm:stochastic-control:problem}{value function} satisfies a
\emph{quasi-variational inequality}\index{quasi-variational inequality}: $\max(\mathcal LV - rV - f,\ \mathcal MV - V) = 0$ (for a
cost-minimisation, with signs reversed), where $\mathcal MV(x) = \sup_\xi\{V(x + \xi) - K - \text{cost}(\xi)\}$ is the value of acting
now: the obstacle depends on the unknown $V$ itself.
\end{definition}

\begin{definition}[No-trade region]\label{def:qm:optimal-stopping-and-impulse-control:notrade}
The \emph{no-trade region}\index{no-trade region} of a trading or hedging policy is the set of states in which it does not
trade; for one exposure it is an interval $(-b, b)$, and a \emph{band policy} trades back to $\pm a$ when the exposure reaches
$\pm b$.
\end{definition}

For the desk, let the exposure move as $\sigma W_t$ between hedges, charge $\gamma X_t^2$ per unit time for carrying it, and
charge $K$ per hedge ticket plus $c$ per unit traded.

\begin{proposition}[The cost of a band policy]\label{prop:qm:optimal-stopping-and-impulse-control:band}
The long-run average cost of the band policy $(a, b)$ is
\[
C(a, b) = \Bigl(K + c(b - a) + \frac{\gamma(b^4 - a^4)}{6\sigma^2}\Bigr)\frac{\sigma^2}{b^2 - a^2}.
\]
With a fixed cost only ($c = 0$) the optimal reset is $a = 0$, and $C(0, b) = K\sigma^2/b^2 + \gamma b^2/6$ is minimised at
\[
b^* = \Bigl(\frac{6K\sigma^2}{\gamma}\Bigr)^{1/4}, \qquad C^* = 2\sqrt{K\sigma^2\gamma/6},
\]
where the fees and the risk charge each contribute half.
\end{proposition}

\begin{proof}
A cycle runs from $\pm a$ to the first exit of $(-b, b)$. By Dynkin's formula (chapter~4), $m(x) = \E_x[\tau]$ and $f(x) =
\E_x\int_0^\tau X_t^2\,dt$ solve $\tfrac12\sigma^2m^{\prime\prime} = -1$ and $\tfrac12\sigma^2f^{\prime\prime} = -x^2$ with zero boundary values: $m(x) = (b^2 -
x^2)/\sigma^2$ and $f(x) = (b^4 - x^4)/(6\sigma^2)$. Each cycle costs $K + c(b - a) + \gamma f(a)$; by the renewal--reward theorem the
average cost is the expected cost per cycle over the expected cycle length. With $c = 0$, $\partial_bC = 0$ gives $b^4 =
6K\sigma^2/\gamma$, and substituting gives $C^*$ with equal parts.
\end{proof}

With $\sigma = 20$ (USD millions per square-root day), $\gamma = \text{USD}~0.5$ per (USD million)$^2$ per day and $K = \text{USD}~300$,
$b^* = 34.6$ million and $C^* = \text{USD}~200$ a day, with a hedge every $b^{*2}/\sigma^2 = 3$ days on average
(\cref{fig:qm:optimal-stopping-and-impulse-control:exposure,fig:qm:optimal-stopping-and-impulse-control:cost}).
Hedging once a day at the close instead costs $K + \gamma\sigma^2/2 = 400$ a day. The cost curve is flat near its minimum and
steep toward small bands: a band of 10 million costs 1\,208 a day.

\begin{omfigure}
\begin{tikzpicture}
\begin{axis}[width=11cm,height=4.8cm,
  xlabel={days},ylabel={exposure (USD m)},
  tick label style={font=\small},label style={font=\small},
  xmin=0,xmax=10,ymin=-42,ymax=42,grid=major,grid style={gray!25},
  table/col sep=comma]
\addplot[omDef,thin] table[x=day,y=x]{figdata/methods/10-optimal-stopping-and-impulse-control/exposure.csv};
\addplot[omThm,dashed,thick,domain=0:10] {34.64};
\addplot[omThm,dashed,thick,domain=0:10] {-34.64};
\end{axis}
\end{tikzpicture}

\omcaption{Ten simulated days of the desk's exposure under the optimal band $\pm 34.6$ million (dashed): client flows move it
freely inside the \omterm{def:qm:optimal-stopping-and-impulse-control:notrade}{no-trade region}, and each touch of the band triggers a hedge back to zero. Data: the chapter's tutorial,
seeded.}\label{fig:qm:optimal-stopping-and-impulse-control:exposure}
\end{omfigure}

\begin{omfigure}
\begin{tikzpicture}
\begin{axis}[width=11cm,height=4.8cm,
  xlabel={band half-width $b$ (USD m)},ylabel={cost (USD a day)},
  tick label style={font=\small},label style={font=\small},
  xmin=8,xmax=90,ymin=0,ymax=1000,grid=major,grid style={gray!25},
  legend columns=2,legend style={font=\footnotesize,at={(0.5,-0.3)},anchor=north,draw=none,fill=none,/tikz/every even column/.append style={column sep=8pt}},
  table/col sep=comma]
\addplot[omDef,very thick] table[x=b,y=cost]{figdata/methods/10-optimal-stopping-and-impulse-control/cost.csv};
\addplot[omThm,dashed,thick] table[x=b,y=daily]{figdata/methods/10-optimal-stopping-and-impulse-control/cost.csv};
\addplot[black,only marks,mark=*,mark size=2pt] coordinates {(34.64,200)};
\legend{band policy $K\sigma^2/b^2 + \gamma b^2/6$,hedge daily at the close}
\end{axis}
\end{tikzpicture}

\omcaption{Average daily cost of the band policy against the band's half-width, for the desk's numbers ($K = 300$,
$\sigma = 20$, $\gamma = 0.5$): minimum USD~200 at $b^* = 34.6$ million (dot), half the cost of hedging once a day. Data:
closed form, checked by simulation in the tutorial.}\label{fig:qm:optimal-stopping-and-impulse-control:cost}
\end{omfigure}

\section{Singular control and reflection}

When the cost is proportional only ($K = 0$), small trades cost little, and the optimal policy trades continuously but only
at the boundary.

\begin{definition}[Singular control, reflected Brownian motion]\label{def:qm:optimal-stopping-and-impulse-control:singular}
A \emph{singular control}\index{singular control} acts through a process of finite variation that increases only on a set
of times of Lebesgue measure zero. \emph{Reflected Brownian motion}\index{reflected Brownian motion} on $[-b, b]$ is $X =
\sigma W + L - U$, with $L$ and $U$ nondecreasing, increasing only when $X = -b$ and $X = b$ respectively: the minimal pushing that
keeps $X$ in the interval.
\end{definition}

\omterm{def:qm:optimal-stopping-and-impulse-control:singular}{Reflected Brownian motion} on $[-b, b]$ has the uniform stationary law, so $\E[X^2] = b^2/3$, and its boundaries push at a total
long-run rate of $\sigma^2/(2b)$ (apply Itô's formula to $X^2$: the drift $\sigma^2$ is balanced by $2b$ times the pushing rate). The
average cost is $c\sigma^2/(2b) + \gamma b^2/3$, minimised at
\[
b^* = \Bigl(\frac{3c\sigma^2}{4\gamma}\Bigr)^{1/3}:
\]
the cube-root law for proportional costs, against the fourth root for fixed costs
(\cref{fig:qm:optimal-stopping-and-impulse-control:laws}). With $c = \text{USD}~25$ per million (a quarter of a basis point),
$b^* = 24.7$ million and the cost is USD~304 a day. With both costs, the optimal band triggers at 42.4 million and trades back
to 12.8 million, not to zero, at USD~418 a day: the fixed fee makes trades large, the proportional cost makes them stop short.
Davis and Norman (1990) solved the portfolio version, a \omterm{def:qm:optimal-stopping-and-impulse-control:notrade}{no-trade region} around the \omterm{def:qm:stochastic-control:fraction}{Merton fraction} of chapter~9.

\begin{omfigure}
\begin{tikzpicture}
\begin{axis}[width=10cm,height=4.8cm,xmode=log,ymode=log,
  xlabel={cost relative to the desk's ($K/300$ or $c/25$)},ylabel={optimal band (USD m)},
  tick label style={font=\small},label style={font=\small},
  xmin=0.1,xmax=10,ymin=9,ymax=70,grid=major,grid style={gray!25},
  ytick={10,20,40,60},yticklabels={10,20,40,60},
  legend columns=2,legend style={font=\footnotesize,at={(0.5,-0.3)},anchor=north,draw=none,fill=none,/tikz/every even column/.append style={column sep=8pt}},
  table/col sep=comma]
\addplot[omDef,very thick] table[x=scale,y=b_fixed]{figdata/methods/10-optimal-stopping-and-impulse-control/laws.csv};
\addplot[omThm,very thick] table[x=scale,y=b_prop]{figdata/methods/10-optimal-stopping-and-impulse-control/laws.csv};
\legend{fixed fee: slope $1/4$,proportional cost: slope $1/3$}
\end{axis}
\end{tikzpicture}

\omcaption{Optimal band against the size of the trading cost, on log scales: with a fixed fee per ticket the band grows like the
fourth root of the fee, with a proportional cost like the cube root. A hundredfold change in cost moves the bands by factors of
3.2 and 4.6. Data: closed forms, the chapter's tutorial.}\label{fig:qm:optimal-stopping-and-impulse-control:laws}
\end{omfigure}

\section{Tutorial: the hedging band}\label{tut:qm:optimal-stopping-and-impulse-control:band}

\begin{tutorial}
\textbf{Goal.} Compute the desk's optimal band, check it by simulating the exposure, find the combined band with fixed and
proportional costs, and solve the take-profit problem by \omterm{def:qm:optimal-stopping-and-impulse-control:fbp}{smooth pasting} and on a grid. \textbf{End state:}
\cref{fig:qm:optimal-stopping-and-impulse-control:exposure,fig:qm:optimal-stopping-and-impulse-control:cost,fig:qm:optimal-stopping-and-impulse-control:laws,fig:qm:optimal-stopping-and-impulse-control:stopping}
and the numbers of the weekend problem.
\begin{tutsteps}
\item \textbf{The band's cost} by renewal--reward, and the two closed-form optima.
\omcode{code/firm/impulse/firm_impulse.py}{26}{50}{The average cost of a band policy, of reflection, and the fourth- and cube-root bands.}\label{lst:qm:optimal-stopping-and-impulse-control:band}
\item \textbf{The stopping threshold} by \omterm{def:qm:optimal-stopping-and-impulse-control:fbp}{smooth pasting}.
\omcode{code/firm/impulse/firm_impulse.py}{95}{100}{The take-profit threshold for a drifting position.}\label{lst:qm:optimal-stopping-and-impulse-control:stop}
\item \textbf{Run} \texttt{problem()}, \texttt{dp\_stopping\_check()} and \texttt{fig\_bands.py}; the simulation of 20\,000 days at
390 checks a day gives USD~197 a day for the optimal band, and 200 with checks four times as frequent.
\end{tutsteps}
\textbf{What to change next.} Make the client flows mean-revert (the exposure an \omterm{def:qm:stochastic-differential-equations:ou}{Ornstein--Uhlenbeck process}) and find how the band
widens; hedge two correlated currencies with one band on their risk-weighted combination.
\end{tutorial}

\section{Build: hedging bands}\label{bld:qm:optimal-stopping-and-impulse-control:impulse}

\begin{build}
\bfield{Purpose} Decide when the firm's hedgers, rebalancers and inventory managers trade: every recurring hedge goes through a
band whose width is computed, not guessed.
\bfield{Interface} \texttt{band\_cost(b, a, sigma, gamma, K, c)}; \texttt{reflect\_cost(b, sigma, gamma, c)};
\texttt{optimal\_fixed\_band}; \texttt{optimal\_proportional\_band}; \texttt{optimal\_band(sigma, gamma, K, c)} returning $(a^*, b^*,
C^*)$; \texttt{simulate\_band}; \texttt{stop\_threshold\_drift(mu, sigma, r, cost)}; \texttt{golden}.
\bfield{Rules} Units stated at the call site (exposure, time, currency); the closed forms are used where they exist and checked by
simulation in the tests; the combined band is searched over both the trigger and the reset.
\bfield{Acceptance tests} \texttt{code/firm/impulse/tests/}: the renewal formula against simulation; the fourth- and cube-root laws;
the numerical optimum reduces to the closed forms when one cost vanishes; \omterm{def:qm:optimal-stopping-and-impulse-control:fbp}{smooth pasting} holds at the stopping threshold.
\bfield{Stretch} Mean-reverting exposures; several currencies with a correlation matrix; the Davis--Norman portfolio \omterm{def:qm:optimal-stopping-and-impulse-control:notrade}{no-trade
region}.
\end{build}

\begin{omsources}
\item J.~L.~Snell, ``Applications of martingale system theorems'', \emph{Transactions of the AMS} 73, 1952.
\item G.~M.~Constantinides and S.~F.~Richard, ``Existence of optimal simple policies for discounted-cost inventory and cash
management in continuous time'', \emph{Operations Research} 26, 1978.
\item M.~H.~A.~Davis and A.~R.~Norman, ``Portfolio selection with transaction costs'', \emph{Mathematics of Operations Research}
15, 1990.
\item A.~E.~Whalley and P.~Wilmott, ``An asymptotic analysis of an optimal hedging model for option pricing with transaction
costs'', \emph{Mathematical Finance} 7, 1997.
\item G.~Peskir and A.~Shiryaev, \emph{Optimal Stopping and Free-Boundary Problems}, Birkhäuser, 2006.
\end{omsources}

\section{Exercises}

\begin{exercise}[$\star$]\label{exo:qm:optimal-stopping-and-impulse-control:1}
With $K = 300$, $\sigma = 20$ and $\gamma = 0.5$, what are the optimal band, its daily cost and the average number of hedges a day?
\end{exercise}

\begin{exercise}[$\star$]\label{exo:qm:optimal-stopping-and-impulse-control:2}
On the two-step tree of \cref{fig:qm:probability-at-speed:tree}, compute the \omterm{def:qm:optimal-stopping-and-impulse-control:snell}{Snell envelope} of $G_n = (S_n - 100)^+$ and say where it
is optimal to stop.
\end{exercise}

\begin{exercise}[$\star$]\label{exo:qm:optimal-stopping-and-impulse-control:3}
Is ``the day the exposure reaches its largest absolute value this month'' a valid time to hedge in the sense of this chapter?
\end{exercise}

\begin{exercise}[$\star\star$]\label{exo:qm:optimal-stopping-and-impulse-control:4}
If the ticket fee doubles, by how much do the optimal band and its cost change? And if the flow volatility doubles?
\end{exercise}

\begin{exercise}[$\star\star$]\label{exo:qm:optimal-stopping-and-impulse-control:5}
With only a proportional cost of USD~25 per million, what are the reflecting band and its cost?
\end{exercise}

\begin{exercise}[$\star\star$]\label{exo:qm:optimal-stopping-and-impulse-control:6}
Derive $b^* = c + 1/\theta$ in \cref{prop:qm:optimal-stopping-and-impulse-control:drift} from \omterm{def:qm:optimal-stopping-and-impulse-control:fbp}{smooth pasting} alone, and evaluate it for
$\mu = 0.5$, $\sigma = 2$, $r = 0.02$, $c = 5$.
\end{exercise}

\begin{exercise}[$\star\star\star$]\label{exo:qm:optimal-stopping-and-impulse-control:7}
\emph{Coding.} Simulate the desk's exposure with the optimal band, checking it 390 and then 1\,560 times a day, and compare the
average cost with the formula. Why is the discretely checked band cheaper?
\end{exercise}

\begin{exercise}[$\star\star\star$]\label{exo:qm:optimal-stopping-and-impulse-control:8}
\emph{Find the flaw.} ``We hedge as soon as the exposure exceeds 10 million: a tighter band means less risk, and less risk is
always better.''
\end{exercise}

\section{Problem: The Hedging Band}

\begin{problem}[{Weekend problem --- when to hedge a drifting exposure}]\label{pb:qm:optimal-stopping-and-impulse-control:1}
A treasury desk's currency exposure $X$ moves with client flows as $\sigma W_t$, $\sigma = 20$ million per square-root day. Carrying it costs
$\gamma X^2$ a day with $\gamma = 0.5$ dollar per (million)$^2$; a hedge ticket costs $K = 300$ dollars; later, trading also costs $c = 25$ dollars
per million.

\medskip\noindent\textbf{Part I --- The fixed fee.}
\begin{enumerate}
\item Write the average cost of the band policy that hedges to zero at $\pm b$.
\item What is the optimal half-width?
\item What does it cost a day?
\item How often does the desk hedge?
\item What does hedging once a day at the close cost, and what is the band's saving?
\end{enumerate}

\medskip\noindent\textbf{Part II --- Scaling.}
\begin{enumerate}[resume]
\item If the fee halves, what are the new band and cost?
\item If the flow volatility doubles?
\item What share of the optimal cost is the risk charge?
\item What does simulation at 390 and at 1\,560 checks a day give?
\item Why does the discretely checked band come out slightly cheaper?
\end{enumerate}

\medskip\noindent\textbf{Part III --- Proportional costs.}
\begin{enumerate}[resume]
\item With the proportional cost only, what are the band and the cost?
\item With both costs, what are the trigger, the reset and the cost?
\item Why does the desk no longer trade back to zero?
\item What distinguishes the \omterm{def:qm:optimal-stopping-and-impulse-control:singular}{singular control} of item~11 from the \omterm{def:qm:optimal-stopping-and-impulse-control:impulse}{impulse control} of item~12?
\item Solve the take-profit problem of \cref{prop:qm:optimal-stopping-and-impulse-control:drift} with the chapter's numbers, in closed form and on a daily grid.
\end{enumerate}

\medskip\noindent\textbf{Part IV --- Judgement.}
\begin{enumerate}[resume]
\item Is a quadratic risk charge a sensible model of the desk's risk?
\item If client flows mean-revert, should the band be wider or narrower?
\item Why should the desk net its currencies before banding them?
\item State the \emph{named result}: the optimal band and its scaling with the fee.
\item In one sentence: why is doing nothing inside a band optimal?
\end{enumerate}
\end{problem}

\section{Interview questions}

\begin{interviewq}[$\star$ \iqroles{trader, bank}]\label{iq:qm:optimal-stopping-and-impulse-control:1}
You must hedge an exposure that drifts with client flows, and each hedge costs a fixed fee. Describe the optimal policy.
\end{interviewq}

\begin{interviewq}[$\star\star$ \iqroles{researcher}]\label{iq:qm:optimal-stopping-and-impulse-control:2}
What is \omterm{def:qm:optimal-stopping-and-impulse-control:fbp}{smooth pasting}, and why should the \omterm{def:qm:stochastic-control:problem}{value function} meet the payoff with the same slope?
\end{interviewq}

\begin{interviewq}[$\star\star$ \iqroles{researcher, trader}]\label{iq:qm:optimal-stopping-and-impulse-control:3}
Why does an optimal hedging band scale with the fourth root of a fixed fee but the cube root of a proportional cost?
\end{interviewq}

\begin{interviewq}[$\star\star$ \iqroles{bank, researcher}]\label{iq:qm:optimal-stopping-and-impulse-control:4}
What is the \omterm{def:qm:optimal-stopping-and-impulse-control:snell}{Snell envelope}, and how does it relate to pricing an American option?
\end{interviewq}

\begin{interviewq}[$\star\star$ \iqroles{developer}]\label{iq:qm:optimal-stopping-and-impulse-control:5}
Implement a band hedger that must not trade back and forth when the exposure sits on the band. What goes into the design?
\end{interviewq}

\begin{interviewq}[$\star\star\star$ \iqroles{researcher}]\label{iq:qm:optimal-stopping-and-impulse-control:6}
For a \omterm{def:qm:brownian-motion:bm}{Brownian motion} with variance $\sigma^2$ started at 0, compute $\E\int_0^\tau X_t^2\,dt$, where $\tau$ is the first exit from $(-b, b)$.
\end{interviewq}
